\exercisesection{Regular Commutative Banach Algebras}

In the exercises below, all algebras considered are over C.

\begin{enumerate}
\def\labelenumi{\arabic{enumi})}
\item
  Let A be a regular commutative Banach algebra and I a closed ideal of A. Then I and A/I are regular commutative Banach algebras.
\item
  Let A be a regular commutative Banach algebra without radical, I an ideal of A, \(\chi\) an element of V(I), and J the set of \(x \in A\) such that Gx has compact support disjoint from \(\{\chi\}\) . Then I + J is the smallest of the ideals R of A such that R ⊃ I and V(R) = \(\{\chi\}\) .
\item
  Let A be a commutative Banach algebra, A' the dual Banach space of A. For every \(x \in A\) , let \(\kappa(x)\) the linear operator in A' transposed to multiplication by x in A. For \(x \in A\) and \(x' \in A'\) , we denote by \(x \star x' = \kappa(x)(x') \in A'\) . A vector subspace of A' is said to be invariant if it is stable under all the \(\kappa(x)\) .
\end{enumerate}

\begin{enumerate}
\def\labelenumi{\alph{enumi})}
\item
  For an element \(x'\) of A' to be proportional to a nonzero character of A, it is necessary and sufficient that \(Cx'\) be invariant.
\item
  The map \(V \mapsto V^{\circ}\) is a bijection from the set of closed ideals of A onto the set of weakly closed invariant vector subspaces of \(A'\) .
\item
  If W is an invariant vector subspace of A', we denote by \(\sigma(\mathrm{W})\) the set of characters belonging to W. If \(W_{1}\) and \(W_{2}\) are weakly closed invariant subspaces, and if W is the weakly closed vector subspace generated by \(W_{1}\) and \(W_{2}\) , one has \(\sigma(\mathrm{W}) = \sigma(\mathrm{W}_{1}) \cup \sigma(\mathrm{W}_{2})\) (Use b).
\item
  Let W be a weakly closed invariant vector subspace of A′. Then \(\sigma(\mathrm{W})\) is the set of \(\chi \in \widehat{A}\) such that the condition \(x * x' = 0\) for all \(x' \in W\) implies \(\chi(x) = 0\) .
\item
  If A admits a unit element, every weakly closed invariant vector subspace of \(A'\) contains a nonzero character.
\item
  In the remainder of this exercise, we assume that A is a regular algebra without radical and with identity. If W is a weakly closed invariant vector subspace of A' and U a neighborhood of \(\sigma(\mathrm{W})\) into \(\mathsf{X}(\mathsf{A})\) then W is contained in the weakly closed vector subspace of A' generated by U.
\end{enumerate}

\begin{enumerate}
\def\labelenumi{\alph{enumi})}
\setcounter{enumi}{6}
\tightlist
\item
  Let \(x' \in A'\) and let W be the weakly closed invariant vector subspace of \(A'\) generated by \(x'\) . We set \(\sigma(x') = \sigma(\mathrm{W})\) . If \(x \in A\) is such that \(\mathcal{G}x = 1\) on a neighborhood of \(\sigma(x')\) , then \(x \star x' = x'\) . (Use \(f\) ).)
\end{enumerate}

\begin{enumerate}
\def\labelenumi{\alph{enumi})}
\setcounter{enumi}{7}
\tightlist
\item
  If \(\sigma(x')\) is the union of two disjoint closed sets \(\sigma_1\) and \(\sigma_2\) , there exists \(x_1', x_2' \in A'\) such that \(\sigma(x_1') = \sigma_1\) , \(\sigma(x_2') = \sigma_2\) and \(x' = x_1' + x_2'\) . (Set \(x_1' = u_1 \star x'\) , \(x_{2}^{\prime}=u_{2}\star x^{\prime}\) , with \(\mathcal{G}u_{1}=1\) (resp. 0) in a neighborhood of \(\sigma_{1}\) (resp. \(\sigma_{2}\) ), and \(\mathcal{G}u_{2}=1\) (resp. 0) in a neighborhood of \(\sigma_{2}\) (resp. \(\sigma_{1}\) ).)
\end{enumerate}

\begin{enumerate}
\def\labelenumi{\arabic{enumi})}
\setcounter{enumi}{3}
\item
  Let A be a commutative regular semisimple Banach algebra satisfying Ditkin's condition, and let I be a closed ideal of A. \(x \in \Upsilon(\mathrm{V}(\mathrm{I}))\) and F the boundary of \(\mathrm{V}(Ax)\) . If \(\mathrm{F} \cap \mathrm{V}(\mathrm{I})\) contains no nonempty perfect set, we have \(x \in I\) (Imitate the reasoning of Prop. 5 of I, p.~94.)
\item
  Let A be a regular commutative Banach algebra without radical. We assume that every closed ideal of A is an intersection of regular maximal ideals. Let \(x \in A\) . Let F be the set of zeros of \(\mathcal{G}'(x)\) Then x is a limit of elements \(u_{n}x\) , where \(u_{n} \in A\) and \(\mathcal{G}(u_{n})\) vanishes in a neighborhood of F. (Let J be the set of \(y \in A\) such that \(\mathcal{G}(y)\) with compact support disjoint from F. We have \(x \in \overline{J}\) . Use prop. 4 of I, p.~91.) In particular, A satisfies Ditkin's condition.
\item
  Let B be the algebra of continuously differentiable complex functions on {[}0, 1{]}, equipped with the norm \(\|f\| = \sup |f| + \sup |f'|\) . Let A be the closed subalgebra of B formed by the functions \(f \in \mathrm{B}\) such that \(f\left(\frac{1}{2}\right) = 0\) . Let I be the closed ideal of A formed by the \(f \in \mathrm{A}\) such that \(f'\left(\frac{1}{2}\right) = 0\) .
\end{enumerate}

\begin{enumerate}
\def\labelenumi{\alph{enumi})}
\item
  The space X(A) is identified with {[}0, 1{]} = \{ \(\frac{1}{2}\) \}, and A is regular.
\item
  The ideal I is a maximal ideal of A that is not regular.
\end{enumerate}

\begin{enumerate}
\def\labelenumi{\arabic{enumi})}
\setcounter{enumi}{6}
\tightlist
\item
  Let R, \(T_{1}\) , \(T_{2}\) be the subsets formed by \(C^{4}\) consisting of the elements \((z_{1}, z_{2}, z_{3}, z_{4}) \in \mathbf{C}^{4}\) such that:
\end{enumerate}

\[
\mathrm{R} \colon z _ {1} z _ {2} = 2, \qquad 1 \leqslant | z _ {1} | \leqslant 2, \qquad z _ {3} = z _ {4} = 0,
\]

\[
\mathrm{T} _ {1} \colon z _ {1} z _ {2} = 2, \qquad | z _ {1} | = 1, \qquad | z _ {3} | \leqslant 1, \qquad z _ {4} = 0,
\]

\[
\mathrm{T} _ {2} \colon z _ {1} z _ {2} = 2, \qquad | z _ {1} | = 2, \qquad | z _ {3} | \leqslant 1, \qquad z _ {4} = z _ {3} ^ {2}.
\]

\begin{enumerate}
\def\labelenumi{\alph{enumi})}
\tightlist
\item
  The set \(X = R \cup T_{1} \cup T_{2}\) is polynomially convex. (Show that X is the set of \((z_{1}, z_{2}, z_{3}, z_{4})\) such that
\end{enumerate}

\[
z _ {1} z _ {2} = 2, \quad | z _ {1} | \leqslant 2, \quad | z _ {2} | \leqslant 2, \quad | z _ {3} | \leqslant 1, \quad z _ {4} (z _ {4} - z _ {3} ^ {2}) = 0,
\]

\[
| z _ {4} z _ {2} ^ {k} | \leqslant 1 \quad \text { and } \quad | z _ {1} ^ {k} (z _ {4} - z _ {3} ^ {2}) | \leqslant 1 \quad \text { for } \quad k = 1, 2, \ldots).
\]

\begin{enumerate}
\def\labelenumi{\alph{enumi})}
\setcounter{enumi}{1}
\item
  Let A be the closed subalgebra of \(\mathcal{C}(\mathrm{X})\) generated by the restrictions to X of polynomial functions on \(\mathbf{C}^4\) . Then X is identified with \(\mathsf{X}(\mathsf{A})\) .
\item
  Let \(\Gamma_1, \Gamma_2 \subset \mathbf{C}\) the circles centered at 0 with radii 1, 2 oriented in the positive direction. Let \(\mu_0, \mu_1, \nu\) the measures on \(\Gamma_1, \Gamma_2, \Gamma_1\) defined by the differential forms \(z^{-1} dz, -z^{-1} dz, z^{-2} dz\) . Let \(\mu = \mu_0 + \mu_1\) , so that \(\mu \otimes \nu\) is a measure on \((\Gamma_1 \cup \Gamma_2) \times \Gamma_1\) . Let B be the set of \(z = (z_1, z_2, z_3, z_4) \in X\) such that \((z_1, z_3) \in (\Gamma_1 \cup \Gamma_2) \times \Gamma_1\) . The map \(\varphi: z \mapsto (z_1, z_3)\) from B into \((\Gamma_{1} \cup \Gamma_{2}) \times \Gamma_{1}\) is a homeomorphism. Let \(\nu\) the measure \(\varphi^{-1}(\mu \otimes \nu)\) . Show that \(\nu\) is orthogonal to A.
\item
  Let g be the function on X that is zero on \(R \cup T_{1}\) and equal to \(z_{3}\) on \(T_{2}\) . Then \(g \notin A\) .
\item
  The function g coincides with an element of A in a neighborhood of every \(z \in X\) .
\end{enumerate}

